\section{The termination classification}\label{sec:classification}

\providecommand{\Cscen}[1]{\mathcal C_{#1}}


\begin{definition}[Double-star]\label{def:doublestar}
A \emph{double-star} is a graph with a distinguished edge $\{b,c\}$, the \emph{central edge}, such that every other vertex is adjacent to exactly one of $b$ and $c$ and to nothing else. Equivalently, a double-star is a tree of diameter at most three. Single edges and stars are double-stars.
\end{definition}

\begin{theorem}[Termination classification]\label{thm:classification}\label{thm:classification-intro}
Let $G$ be a finite simple graph without isolated vertices, with binary observations. The following are equivalent.
\begin{enumerate}[label=(\alph*),nosep]
\item $\INW_t(G)=\Cscen{G}$ for some finite $t$.
\item $\IAI_t(G)=\Cscen{G}$ for some finite $t$.
\item $\Iexp_t(G)=\Cscen{G}$ for some finite $t$.
\item Every connected component of $G$ is a double-star.
\end{enumerate}
When they hold,
\[
\INW_2(G)=\IAI_2(G)=\Iexp_2(G)=\Cscen{G}.
\]
When they fail, for every $t\ge1$ there is a strictly positive rational law
\[
P\in\Iexp_t(G)\setminus\Cscen{G}\subseteq\IAI_t(G)\setminus\Cscen{G}\subseteq\INW_t(G)\setminus\Cscen{G}.
\]
\end{theorem}

The graph and the alphabets stay fixed as $t$ grows; only the law $P$ varies. The proof combines Theorem~\ref{thm:doublestar} with Theorems~\ref{thm:cycle} and \ref{thm:fivepath} through the exhaustion lemma below.

\subsection{Structural lemmas}

In the inflated graph every source copy is a root and every copied observation is a sink, and as a result the AI and expressible prescriptions coincide; Appendix~\ref{app:rootsink} proves the equivalence.

\begin{definition}[Shared-parent graph and AI sets]\label{def:sharedparent}
For a set $W$ of copied observations, the \emph{shared-parent graph} on $W$ joins two members when they read a common source copy. Call $W$ an \emph{AI set} if it is a union of pairwise ancestrally independent injectable sets.
\end{definition}

\begin{lemma}[Root-sink lemma]\label{lem:rootsink}
Let $G$ be a pair-source scenario and $t\ge1$.
\begin{enumerate}[label=(\roman*),nosep]
\item A set $W$ of copied observations is an AI set if and only if every connected component of its shared-parent graph is injectable.
\item A table satisfies every expressible prescription of Definition~\ref{def:gexp} if and only if it satisfies the injectable and ancestral-independence prescriptions of Definition~\ref{def:gAI}.
\end{enumerate}
In particular $\Iexp_t(G)=\IAI_t(G)$ for every $t$.
\end{lemma}


Lemma~\ref{lem:rootsink} identifies the two stronger tests. It says nothing about the Navascués--Wolfe test and imposes no causal model on the table as a whole.

\paragraph{Source-disjoint blocks.}

\begin{lemma}[Source-disjoint independence]\label{lem:sourcedisjoint}
Let $t\ge2$ and $P\in\INW_t(G)$. If $U,W\subseteq V$ satisfy $E(U)\cap E(W)=\varnothing$, then
\[
P_{U\cup W}=P_U\otimes P_W.
\]
\end{lemma}

\begin{proof}
Let $\Gamma$ be a witness. Let $\sigma$ be the element of $\prod_eS_t$ that transposes the copy indices $1$ and $2$ on every edge of $E(W)$ and fixes all other indices. For $u\in U$ every edge at $u$ lies outside $E(W)$, so $\sigma$ fixes $O_u^{\,1}$; for $w\in W$ every edge at $w$ lies in $E(W)$, so $\sigma$ sends $O_w^{\,1}$ to $O_w^{\,2}$. Invariance of $\Gamma$ gives
\[
\Law\bigl((O_u^{\,1})_{u\in U},(O_w^{\,1})_{w\in W}\bigr)=\Law\bigl((O_u^{\,1})_{u\in U},(O_w^{\,2})_{w\in W}\bigr).
\]
The left side is the marginal of row $1$, namely $P_{U\cup W}$. The right side reads $U$-coordinates from row $1$ and $W$-coordinates from row $2$, and the two rows are independent with law $P$ each, so it is $P_U\otimes P_W$.
\end{proof}

Two vertices joined by an edge have that edge in both incidence sets, so the hypothesis says exactly that $U$ and $W$ are disjoint and no edge joins them. Iterating the lemma with $U$ a single block and $W$ the union of the others gives mutual independence of any family of pairwise source-disjoint blocks. In particular the observed law of a disconnected graph factors over its components from order two on.

\paragraph{Induced subgraphs.}

\begin{lemma}[Induced-subgraph transport]\label{lem:transport}
Let $H$ be an induced subgraph of $G$ in which no vertex is isolated, and let $\mathcal J\in\{\INW,\IAI,\Iexp\}$.
\begin{enumerate}[label=(\roman*),nosep]
\item If $Q\in\Cscen{G}$ then $Q_{V(H)}\in\Cscen{H}$.
\item If $P_H\in\mathcal J_t(H)$ then $P:=P_H\otimes\Bern(1/2)^{\otimes(V\setminus V(H))}\in\mathcal J_t(G)$.
\item $d_{\TV}(P,\Cscen{G})\ge d_{\TV}(P_H,\Cscen{H})$. In particular, if $P_H\notin\Cscen{H}$ then $P\notin\Cscen{G}$.
\end{enumerate}
\end{lemma}

The proof is in Appendix~\ref{app:transport}.

Part (ii) preserves strict positivity and rationality of the target, since the appended factor is the fair-bit law.

\subsection{Reduction to cycles and paths}

\begin{lemma}[Exhaustion]\label{lem:exhaustion}
Let $K$ be a finite connected graph that has at least one edge and is not a double-star. Then $K$ has an induced subgraph isomorphic to a cycle $C_m$ with $m\ge3$ or to the path $P_5$.
\end{lemma}

\begin{proof}
Suppose first that $K$ contains a cycle, and let $C$ be one of minimum length. A chord of $C$ would split it into two shorter cycles, so $C$ has no chords and is induced. Its length is at least $3$.

Suppose next that $K$ is a tree. We claim that a tree of diameter at most three is a double-star. If the diameter is $1$ then $K$ is a single edge. If it is $2$ then $K$ is a star. Let the diameter be $3$ and fix a geodesic $u_0u_1u_2u_3$. Let $w$ be a vertex other than $u_1$ and $u_2$, and consider the unique path from $w$ to $u_1$. If it avoids $u_2$, then it also avoids $u_3$, since the path from $u_3$ to $u_1$ passes through $u_2$; so concatenating with $u_1u_2u_3$ gives a path and $d(w,u_3)=d(w,u_1)+2\le3$, whence $d(w,u_1)=1$. If it contains $u_2$, then the path from $w$ to $u_2$ avoids $u_1$ and hence $u_0$, and the same argument with $u_0$ gives $d(w,u_2)=1$. Every vertex is therefore adjacent to $u_1$ or $u_2$, and $K$ is a double-star with central edge $u_1u_2$.

A tree that is not a double-star therefore has diameter at least $4$. Take a geodesic of length at least $4$ and let $v_0,\dots,v_4$ be five consecutive vertices on it. Any edge of $K$ among them other than the four geodesic edges would create a cycle in a tree, so the induced subgraph on $\{v_0,\dots,v_4\}$ is $P_5$.
\end{proof}

Fritz classified the graph scenarios in which every correlation is classical as the disjoint unions of stars~\cite[Theorem 3.8]{Fritz2012}: every other connected graph has an induced triangle, square or four-vertex path~\cite[Lemma 3.9]{Fritz2012}, and a correlation on that induced subgraph extends to the whole graph by deterministic outcomes at the remaining vertices. Lemma~\ref{lem:exhaustion} and the proof below follow that argument, with double-stars in place of stars, cycles and the five-vertex path as the forbidden induced subgraphs, and inflation witnesses in place of correlations. Fritz also framed scenarios with arbitrary sources as hypergraphs~\cite[Section 3]{Fritz2012}.

\begin{proof}[Proof of Theorem~\ref{thm:classification}]
Assume (d). Each component is a double-star with binary observations, so Theorem~\ref{thm:doublestar} and Corollary~\ref{cor:disjointunions} give $\INW_2(G)=\Cscen{G}$. With Proposition~\ref{prop:gsound},
\[
\Cscen{G}\subseteq\Iexp_2(G)\subseteq\IAI_2(G)\subseteq\INW_2(G)=\Cscen{G},
\]
so all four sets coincide and (a), (b), (c) hold at $t=2$.

Assume (d) fails, and fix $t\ge1$. Some component $K$ of $G$ is not a double-star, and $K$ has an edge because $G$ has no isolated vertices. By Lemma~\ref{lem:exhaustion}, $K$ has an induced subgraph $H$ isomorphic to $C_m$ for some $m\ge3$ or to $P_5$; since $K$ is induced in $G$, so is $H$, and $H$ has no isolated vertices. Theorem~\ref{thm:cycle} in the first case and Theorem~\ref{thm:fivepath} in the second supply a strictly positive rational law
\[
P_H\in\Iexp_t(H)\setminus\Cscen{H}.
\]
Let $P=P_H\otimes\Bern(1/2)^{\otimes(V\setminus V(H))}$. By Lemma~\ref{lem:transport}(ii), $P\in\Iexp_t(G)$, and by Lemma~\ref{lem:transport}(iii), $P\notin\Cscen{G}$. It is strictly positive and rational because $P_H$ is. By Proposition~\ref{prop:gsound} the same law lies in $\IAI_t(G)$ and in $\INW_t(G)$, so none of the three tests equals $\Cscen{G}$ at order $t$. As $t$ was arbitrary, (a), (b) and (c) all fail.

This proves the equivalence, the order-two bound and the witness assertion.
\end{proof}

For $m=3$ the witness of Theorem~\ref{thm:cycle} is a second nonterminating family for the triangle, alongside the defect construction of Theorem~\ref{thm:main}.

\begin{remark}[Larger observed alphabets]\label{rem:alphabets}
Fix finite alphabets $\mathcal O_v$ with $|\mathcal O_v|\ge2$. The classification is unchanged, with a different order in the terminating case. If some component is not a double-star, regard the binary law $P$ of Theorem~\ref{thm:classification} as a law on $\prod_v\mathcal O_v$ supported on $\{0,1\}^V$. A compatible realization of such a law has responses that leave $\{0,1\}$ with probability zero, hence is a binary realization, so the law remains incompatible; and its witness, read in the larger alphabets, satisfies the same prescriptions. The transferred law is supported on the binary sub-alphabet, so the strict positivity in Theorem~\ref{thm:classification} is a statement about binary observations. If every component is a double-star, Theorem~\ref{thm:doublestar} gives termination at
\[
T=\max\bigl(\{2\}\cup\{|\mathcal O_v|:v\text{ a leaf}\}\bigr),
\]
the leaves being taken with respect to a choice of central edge in each component; the bound does not involve the alphabets at the central vertices.
\end{remark}

\begin{remark}[Scenarios that are not pair-source]\label{rem:hypergraph}
Let three observers read one common source. That source can carry a sample of the whole observed triple, so every law on the three outcomes is compatible~\cite[Proposition 3.7]{Fritz2012} and the order-one test characterizes the compatible set. Its observer adjacency graph is the same as that of the pair-source triangle, whose hierarchy does not terminate. Thus the adjacency graph alone does not decide termination when a source may be read by three or more observers.
\end{remark}

\subsection{Finite reconstruction on double-stars}\label{sec:doublestar}

\providecommand{\Cscen}[1]{\mathcal C_{#1}}
\providecommand{\supp}{\operatorname{supp}}

Fix a double-star $G$ in the sense of Definition~\ref{def:doublestar}: two central vertices $b,c$ joined by the central edge $y=\{b,c\}$, left leaves $\ell_1,\dots,\ell_p$ joined to $b$ by edges $e_1,\dots,e_p$, and right leaves $r_1,\dots,r_q$ joined to $c$ by edges $f_1,\dots,f_q$. Observed alphabets $\mathcal O_v$ are finite and arbitrary. Put
\[
T=\max\bigl(\{2\}\cup\{|\mathcal O_{\ell_s}|\}_{s\le p}\cup\{|\mathcal O_{r_j}|\}_{j\le q}\bigr).
\]

In the order-$T$ inflation the leaves have one index each, and the centres have one index per incident edge. Write
\[
L_s^{\,m}=O_{\ell_s}^{\,m},\qquad R_j^{\,m}=O_{r_j}^{\,m},\qquad
B^{\,\mathbf i,n}=O_b^{\,(\mathbf i,n)},\qquad C^{\,n,\mathbf k}=O_c^{\,(n,\mathbf k)},
\]
where $m,n\in[T]$, $\mathbf i\in[T]^p$ records the index on each $e_s$, and $\mathbf k\in[T]^q$ the index on each $f_j$. Row $m$ consists of $L_s^{\,m}$, $R_j^{\,m}$, $B^{\,m,m}$ and $C^{\,m,m}$, with $\mathbf i$ and $\mathbf k$ constant equal to $m$.

For a law $P$ on $\prod_v\mathcal O_v$ write $P(\beta,\gamma,\xi,\zeta)$ for the probability that $O_b=\beta$, $O_c=\gamma$, $O_{\ell_s}=\xi_s$ for all $s$ and $O_{r_j}=\zeta_j$ for all $j$.

\begin{theorem}[Double-star reconstruction]\label{thm:doublestar}
For every double-star $G$ with finite observed alphabets,
\[
\INW_T(G)=\Cscen{G}.
\]
\end{theorem}

\begin{proof}
The inclusion $\Cscen{G}\subseteq\INW_T(G)$ is Proposition~\ref{prop:gsound}. Let $P\in\INW_T(G)$ and let $\Gamma$ be a witness at order $T$. We recover a model for $P$ in four steps.

\medskip
\emph{(i) The leaves and all their copies are independent.}
Distinct leaves have disjoint incidence sets, each consisting of a single edge, so Lemma~\ref{lem:sourcedisjoint}, available because $T\ge2$, applies to one leaf against the union of the others and, iterated, gives
\begin{equation}\label{eq:leafprod}
P_{\mathrm{leaf}}(\xi,\zeta):=P\bigl(O_{\ell_s}=\xi_s\ \forall s,\ O_{r_j}=\zeta_j\ \forall j\bigr)=\prod_{s}P_{\ell_s}(\xi_s)\prod_{j}P_{r_j}(\zeta_j).
\end{equation}
Under $\Gamma$ the $T$ rows are independent with law $P$ each, and the leaf entries of a single row are independent by \eqref{eq:leafprod}. Every copied leaf observation occurs in exactly one row, so the whole family $\{L_s^{\,m},R_j^{\,m}\}$ is independent under $\Gamma$, with $L_s^{\,m}\sim P_{\ell_s}$ and $R_j^{\,m}\sim P_{r_j}$.

\medskip
\emph{(ii) A positive conditioning event.}
For each left leaf choose $x_s^{\,1},\dots,x_s^{\,T}\in\mathcal O_{\ell_s}$ so that every symbol of $\supp P_{\ell_s}$ occurs in the list, which is possible because $T\ge|\mathcal O_{\ell_s}|$; fill the remaining positions with repetitions of a supported symbol. Choose $z_j^{\,1},\dots,z_j^{\,T}$ for each right leaf in the same way. Let
\[
E=\bigl\{L_s^{\,m}=x_s^{\,m}\ \forall s,m\bigr\}\cap\bigl\{R_j^{\,m}=z_j^{\,m}\ \forall j,m\bigr\}.
\]
By (i),
\begin{equation}\label{eq:QE}
\Gamma(E)=\prod_{s,m}P_{\ell_s}(x_s^{\,m})\prod_{j,m}P_{r_j}(z_j^{\,m})>0.
\end{equation}

\medskip
\emph{(iii) The conditional marginal of the two centres.}
Fix $\mathbf i\in[T]^p$ and $\mathbf k\in[T]^q$. Choose permutations $\sigma_s\in S_T$ of the copy indices of $e_s$ with $\sigma_s(1)=i_s$, permutations $\tau_j\in S_T$ of those of $f_j$ with $\tau_j(1)=k_j$, and the identity on the central edge $y$. The resulting element of the symmetry group carries row $m$ to
\[
\Bigl(B^{\,(\sigma_s(m))_s,\,m},\ C^{\,m,(\tau_j(m))_j},\ \bigl(L_s^{\,\sigma_s(m)}\bigr)_s,\ \bigl(R_j^{\,\tau_j(m)}\bigr)_j\Bigr),
\]
which for $m=1$ is $\bigl(B^{\,\mathbf i,1},C^{\,1,\mathbf k},(L_s^{\,i_s})_s,(R_j^{\,k_j})_j\bigr)$. By invariance the joint law of these $T$ images equals the joint law of the $T$ rows, namely $P^{\otimes T}$. As $m$ runs over $[T]$ the index $\sigma_s(m)$ runs over $[T]$, so every copied leaf observation occurs exactly once among the images and the event $E$ is determined by them. Marginalizing the images of rows $2,\dots,T$ over their two centre entries and using \eqref{eq:leafprod},
\[
\Gamma\bigl(B^{\,\mathbf i,1}=\beta,\ C^{\,1,\mathbf k}=\gamma,\ E\bigr)
=P\bigl(\beta,\gamma,x^{\mathbf i},z^{\mathbf k}\bigr)
\prod_s\prod_{m\ne i_s}P_{\ell_s}(x_s^{\,m})\prod_j\prod_{m\ne k_j}P_{r_j}(z_j^{\,m}),
\]
where $x^{\mathbf i}=(x_s^{\,i_s})_s$ and $z^{\mathbf k}=(z_j^{\,k_j})_j$. Dividing by \eqref{eq:QE} and using \eqref{eq:leafprod} once more,
\begin{equation}\label{eq:condmarg}
\Gamma\bigl(B^{\,\mathbf i,1}=\beta,\ C^{\,1,\mathbf k}=\gamma\mid E\bigr)
=P\bigl(O_b=\beta,\ O_c=\gamma\ \bigm|\ O_{\ell_s}=x_s^{\,i_s}\ \forall s,\ O_{r_j}=z_j^{\,k_j}\ \forall j\bigr),
\end{equation}
the conditioning event on the right having positive probability whenever every $x_s^{\,i_s}$ and $z_j^{\,k_j}$ is supported. This identity uses only the symmetry and the diagonal law of the witness.

\medskip
\emph{(iv) Two response tables and fresh sources.}
For $\xi\in\prod_s\supp P_{\ell_s}$ pick positions $\mathbf i(\xi)$ with $x_s^{\,i_s(\xi)}=\xi_s$, and for $\zeta\in\prod_j\supp P_{r_j}$ pick $\mathbf k(\zeta)$ with $z_j^{\,k_j(\zeta)}=\zeta_j$. Under the conditional law $\Gamma(\cdot\mid E)$ define the random maps
\[
F(\xi)=B^{\,\mathbf i(\xi),1},\qquad G(\zeta)=C^{\,1,\mathbf k(\zeta)},
\]
with arbitrary fixed values on unsupported arguments. Both are functions of the same table, so the pair $(F,G)$ has a joint law $\mu$, and \eqref{eq:condmarg} reads
\begin{equation}\label{eq:mutable}
\mu\bigl(F(\xi)=\beta,\ G(\zeta)=\gamma\bigr)=P\bigl(O_b=\beta,O_c=\gamma\mid\xi,\zeta\bigr)
\end{equation}
for all supported $\xi,\zeta$.

Now take fresh independent sources: let the central source be $Y^*=(F,G)\sim\mu$, let the source of $e_s$ be $X_s^*\sim P_{\ell_s}$ and the source of $f_j$ be $Z_j^*\sim P_{r_j}$, all mutually independent. Respond by
\[
O_{\ell_s}=X_s^*,\qquad O_{r_j}=Z_j^*,\qquad O_b=F(X_1^*,\dots,X_p^*),\qquad O_c=G(Z_1^*,\dots,Z_q^*).
\]
Each response reads only the sources of its incident edges: $b$ reads $X_1^*,\dots,X_p^*$ and $Y^*$, and $c$ reads $Y^*$ and $Z_1^*,\dots,Z_q^*$. This is a model for the scenario $G$, and by \eqref{eq:mutable} and \eqref{eq:leafprod} its observed law is
\[
\prod_sP_{\ell_s}(\xi_s)\prod_jP_{r_j}(\zeta_j)\,\mu\bigl(F(\xi)=\beta,G(\zeta)=\gamma\bigr)
=P_{\mathrm{leaf}}(\xi,\zeta)\,P(\beta,\gamma\mid\xi,\zeta)=P(\beta,\gamma,\xi,\zeta).
\]
The model never samples unsupported leaf values, so the arbitrary choices there do not matter. Hence $P\in\Cscen{G}$.
\end{proof}

A central source with
\(
|\mathcal O_b|^{\prod_s|\supp P_{\ell_s}|}\,|\mathcal O_c|^{\prod_j|\supp P_{r_j}|}
\)
values suffices: a left response table is a map on the product of the left leaf supports, and likewise on the right. Empty products equal one.

\begin{corollary}[Binary observations]\label{cor:doublestarbinary}
For a double-star with binary observations, $\INW_2(G)=\Cscen{G}$.
\end{corollary}

\begin{proof}
All leaf alphabets have two elements, so $T=2$.
\end{proof}

\begin{corollary}[Four-observer path]\label{cor:path}
Consider $A(X)$, $B(X,Y)$, $C(Y,Z)$, $D(Z)$, the path $P_4$ with $A$ and $D$ at the ends. Then $\INW_T(P_4)=\Cscen{P_4}$ with $T=\max\{2,|\mathcal O_A|,|\mathcal O_D|\}$. For binary endpoints $T=2$, whatever the finite alphabets of $B$ and $C$.
\end{corollary}

\begin{proof}
$P_4$ is the double-star with central edge $\{B,C\}$, one left leaf $A$ and one right leaf $D$; apply Theorem~\ref{thm:doublestar}, whose bound involves only the leaf alphabets.
\end{proof}
Fritz~\cite[Theorem~2.4]{Fritz2012} identifies classical correlations on $P_4$ with Bell-local conditional distributions, with the endpoint observations serving as settings. The four-observer path is the four-on-line scenario of \cite[Figure~8]{NW2020}. Navascu\'es and Wolfe list it with the triangle among the scenarios whose compatible set is cut out by inequality constraints, show that order two does not suffice for the triangle, and leave the four-on-line case open. The corollary settles it at order two for binary endpoints.

\begin{corollary}[Disjoint unions]\label{cor:disjointunions}
Let $G$ be a disjoint union of double-stars $G_1,\dots,G_n$ with orders $T_1,\dots,T_n$ as above, and let $T=\max\{2,T_1,\dots,T_n\}$. Then $\INW_T(G)=\Cscen{G}$.
\end{corollary}

\begin{proof}
Let $P\in\INW_T(G)$ with witness $\Gamma$. Distinct components are source-disjoint, so Lemma~\ref{lem:sourcedisjoint}, iterated, gives $P=\bigotimes_aP_{V(G_a)}$. The restriction of $\Gamma$ to the copied observations at the vertices of $G_a$ is a table for the order-$T$ inflation of $G_a$: it is invariant under permutations of the copy indices of the edges of $G_a$, and its diagonal law is the corresponding marginal of $P^{\otimes T}$, namely $\bigl(P_{V(G_a)}\bigr)^{\otimes T}$. Hence $P_{V(G_a)}\in\INW_T(G_a)\subseteq\INW_{T_a}(G_a)$, the inclusion because discarding the copy indices outside $[T_a]$ leaves a witness at order $T_a$. Theorem~\ref{thm:doublestar} now gives $P_{V(G_a)}\in\Cscen{G_a}$. Realize each component by its own model on its own independent source family; the joint observed law is the product, which is $P$.
\end{proof}

A star is the case $q=0$, with $c$ one of its leaves. There the bound of Theorem~\ref{thm:doublestar} is not sharp: order two suffices with arbitrary finite alphabets \cite[Section 4.1]{NW2020}, by taking independent sources equal to the leaf values and letting the centre respond with its conditional law given them. Theorem~\ref{thm:doublestar} adds the central edge, across which the two centres share a source, so the leaf values on the two sides have to be assembled into one pair of response tables.

\providecommand{\Cscen}[1]{\mathcal C_{#1}}
\providecommand{\Csq}{\mathcal C_{\square}}
\providecommand{\Ccyc}{\mathcal C_{C_m}}
\providecommand{\Real}{\operatorname{Re}}
\providecommand{\sgn}{\operatorname{sgn}}

\section{Cycles}\label{sec:cycles}

Section~\ref{sec:pairsource} writes the three finite tests for an arbitrary pair-source scenario. For every cycle and every order we construct a strictly positive law that the strongest of them accepts and that no model produces, then measure how far the accepted laws can sit from compatibility.

Throughout, the $m$-cycle $C_m$, $m\ge3$, has observed variables $O_0,\dots,O_{m-1}$ and sources $S_0,\dots,S_{m-1}$, indices read modulo $m$, with $O_v$ a child of $S_{v-1}$ and $S_v$; the copied observations of Definition~\ref{def:ginflation} are written $O_v^{ij}$, children of $S_{v-1,i}$ and $S_{v,j}$. The triangle is $C_3$, with $\Ctri=\Cscen{C_3}$, read as $A=O_0$, $B=O_1$, $C=O_2$ and $X=S_0$, $Y=S_1$, $Z=S_2$; the square is $C_4$, written $\square$, read as $A(X,W)$, $B(X,Y)$, $C(Y,Z)$, $D(Z,W)$ with $X=S_0$, $Y=S_1$, $Z=S_2$, $W=S_3$.

\subsection{Parity witnesses on every cycle}\label{sec:cycleparity}

\begin{lemma}[Auxiliary parity density]\label{lem:parityaux}
For $N\ge1$ and $q\ge0$ let
\begin{equation}\label{eq:Hdensity}
H_{N,q}(s)=2^{-N}\sum_{\substack{S\subseteq[N]\\|S|\text{ even}}}(-q)^{|S|/2}\prod_{w\in S}s_w,
\qquad s\in\{-1,1\}^N .
\end{equation}
Then $H_{N,q}$ is normalized, its characters are
\begin{equation}\label{eq:Hmoments}
\E_{H_{N,q}}\Bigl[\prod_{w\in S}s_w\Bigr]=
\begin{cases}
0,&|S|\text{ odd},\\
(-q)^{|S|/2},&|S|\text{ even},
\end{cases}
\end{equation}
and if $N^2q\le1/4$ then $H_{N,q}(s)\ge\tfrac23\,2^{-N}>0$ for every $s$.
\end{lemma}

\begin{proof}
Normalization and \eqref{eq:Hmoments} are orthogonality of the characters of $\{-1,1\}^N$ under the uniform measure. For positivity, $\binom N{2j}\le N^{2j}$, so
\[
\bigl|2^NH_{N,q}(s)-1\bigr|\le\sum_{j\ge1}\binom N{2j}q^j\le\sum_{j\ge1}(N^2q)^j\le\sum_{j\ge1}4^{-j}=\tfrac13 . \qedhere
\]
\end{proof}

For $F\subseteq\{0,\dots,m-1\}$ let $\partial F=\{w:|F\cap\{w,w+1\}|=1\}$ be the set of edges of $C_m$ incident to an odd number of vertices of $F$. Every vertex meets two edges, so $|\partial F|$ is even.

\begin{definition}[Cycle parity family]\label{def:cycletarget}
For $m\ge3$ and $q\ge0$ let $P_{m,q}$ be the signed measure on $\{-1,1\}^m$ with density $2^{-m}\sum_F(-q)^{|\partial F|/2}\prod_{v\in F}o_v$ against the counting measure, so that it has total mass $1$ and
\begin{equation}\label{eq:cycletarget}
\E_{P_{m,q}}\Bigl[\prod_{v\in F}O_v\Bigr]=(-q)^{|\partial F|/2}\qquad\text{for every }F .
\end{equation}
\end{definition}

If $m^2q\le1/4$ then $P_{m,q}$ is a probability law, being the law of $(\sigma_{v-1}\sigma_v)_{v=0}^{m-1}$ for $\sigma\sim H_{m,q}$. The exact ranges in the two cases we use are wider.

Taking $F$ the whole vertex set gives $\E[\prod_vO_v]=1$, so wherever $P_{m,q}$ is a probability law the product of all $m$ signs is $1$ almost surely; taking $F$ a contiguous arc gives mean $-q$. For $m=3$ the family is the parity-even law with observed means $x,y,z$,
\begin{equation}\label{eq:paritylaw}
\Pi(x,y,z)(\alpha,\beta,\gamma)=\tfrac14\bigl(1+x\alpha+y\beta+z\gamma\bigr)\ \text{ on }\alpha\beta\gamma=1,\qquad P_{3,q}=\Pi(-q,-q,-q),
\end{equation}
which is a probability law exactly for $0\le q\le1/3$; and for $m=4$ it is the square parity family $P_q:=P_{4,q}$ with atoms
\begin{equation}\label{eq:squaretarget}
\tfrac18(1-6q+q^2)\ \text{ at }0000,\qquad
\tfrac18(1-q^2)\ \text{ at }0011,0110,1001,1100,\qquad
\tfrac18(1+q)^2\ \text{ at }0101,1010,1111,
\end{equation}
zero on odd parity, in the convention $\text{false}=0\leftrightarrow+1$. It is a law exactly for $0\le q\le3-2\sqrt2$.

Fix $t\ge1$, put $N=mt$ and draw auxiliary signs $s=(s_{v,i})_{v\in\{0,\dots,m-1\},\,i\in[t]}$ from $H_{N,q}$. Define every copied observation by
\begin{equation}\label{eq:cyclewitness}
O_v^{ij}=s_{v-1,i}\,s_{v,j}.
\end{equation}
For a set $F$ of copied observations let $\partial F$ denote the set of copied sources occurring an odd number of times in $\prod_{O\in F}O$; each copied observation contributes two incidences, so $|\partial F|$ is even, and \eqref{eq:Hmoments} gives
\begin{equation}\label{eq:masterchar}
\E\Bigl[\prod_{O\in F}O\Bigr]=(-q)^{|\partial F|/2}.
\end{equation}

\begin{lemma}[Boundaries of prescribed characters]\label{lem:boundary}
Let $F_1,\dots,F_r$ be pairwise ancestrally independent injectable sets of copied observations of $C_m$ and let $U_i\subseteq F_i$. Then each $\partial U_i$ has even cardinality and contains at most one copy of each source, so it is empty or meets at least two source families; the sets $\partial U_1,\dots,\partial U_r$ are pairwise disjoint; and $\partial(U_1\cup\dots\cup U_r)=\bigsqcup_i\partial U_i$, which again is empty or meets at least two source families. Moreover $\partial U_i$ is the image, under erasure of copy indices, of the boundary of the corresponding original observed set.
\end{lemma}

\begin{proof}
An injectable set contains at most one copied observation over each original vertex and its ancestry contains at most one copy of each original source, so $\partial U_i$ contains at most one coordinate from each family; being of even cardinality it is empty or has at least two elements in distinct families. Erasing copy indices carries the incidence count of $U_i$ to that of its image, which gives the last claim. Ancestrally independent sets have disjoint ancestries, hence disjoint coordinate sets and disjoint boundaries, and the boundary of a union of sets with disjoint ancestries is the disjoint union of the boundaries. If that union is nonempty then some $\partial U_i$ is nonempty and already meets two families.
\end{proof}

\begin{lemma}[The parity witness passes every prescription]\label{lem:cyclewitness}
Let $m\ge3$, $t\ge1$ and $0<q\le1/(4m^2t^2)$. The law of \eqref{eq:cyclewitness} is a strictly positive witness placing $P_{m,q}$ in $\IAI_t(C_m)=\Iexp_t(C_m)$.
\end{lemma}

\begin{proof}
Positivity of $H_{N,q}$ holds since $N^2q=m^2t^2q\le1/4$. The density \eqref{eq:Hdensity} is symmetric in the $N$ coordinates, so the witness is invariant under independent permutations of the copy indices of each source. By Lemma~\ref{lem:boundary} an injectable set $F$ has $\partial F$ equal to the boundary of its image, so \eqref{eq:masterchar} and \eqref{eq:cycletarget} give it the prescribed marginal of $P_{m,q}$. For pairwise ancestrally independent injectable $F_1,\dots,F_r$ and any $U_i\subseteq F_i$, Lemma~\ref{lem:boundary} and \eqref{eq:masterchar} give
\[
\E\Bigl[\prod_i\prod_{O\in U_i}O\Bigr]=(-q)^{\frac12\sum_i|\partial U_i|}=\prod_i\E\Bigl[\prod_{O\in U_i}O\Bigr],
\]
and characters determine a law on a finite sign cube, so the blocks are independent with the prescribed marginals. The diagonal rows are pairwise ancestrally independent injectable copied cycles, so the diagonal law is $P_{m,q}^{\otimes t}$. Lemma~\ref{lem:rootsink} supplies the recursive prescriptions.
\end{proof}

The witness is a law on the copied observations, not the output of an inflated model: under $H_{N,q}$ the auxiliary signs are correlated across families.

\begin{lemma}[Exact parity rigidity]\label{lem:parityrigidity}
Let $R\in\Ctri$ with $ABC=1$ almost surely. Then there are $\sigma,\tau,\upsilon\in[-1,1]$ with $(\E A,\E B,\E C)=(\sigma\tau,\sigma\upsilon,\tau\upsilon)$. In particular $\E A\cdot\E B\cdot\E C\ge0$.
\end{lemma}

\begin{proof}
Absorb the response seeds into incident sources, so that $A=A(X,Z)$, $B=B(X,Y)$, $C=C(Z,Y)$ are deterministic signs. For $\mu_Y$-almost every $y$ the identity $A(X,Z)B(X,y)C(Z,y)=1$ holds for $\mu_X\otimes\mu_Z$-almost every $(X,Z)$; fix such a $y_0$ and put $S(X)=B(X,y_0)$, $T(Z)=C(Z,y_0)$, so that $A=ST$ almost surely. Substituting in $ABC=1$ gives $\bigl(S(X)B(X,Y)\bigr)\bigl(T(Z)C(Z,Y)\bigr)=1$ almost surely. Write $B'=SB$ and $C'=TC$. Conditionally on $Y$ these are independent signs, and they agree almost surely, so each equals a function $U(Y)$ almost surely. Hence $A=ST$, $B=SU$, $C=TU$ almost surely, and independence of the three sources gives the means with $\sigma=\E S$, $\tau=\E T$, $\upsilon=\E U$. Their product is $(\sigma\tau\upsilon)^2$.
\end{proof}

Rigidity of the parity-perfect (parity token-counting) triangle correlations is prior work: the exact statement is due to Boreiri, Girardin, Ulu, Lipka-Bartosik, Brunner and Sekatski \cite{BoreiriEtAl2023}, and Boreiri, Ulu, Brunner and Sekatski \cite{BoreiriEtAl2025} give a noise-robust form (a compatible law with parity error $\epsilon$ is within $3\epsilon$ of a parity-perfect compatible law); Lemma~\ref{lem:quantrigidity} is a moment form of that robustness. Renou and Beigi \cite{RenouBeigi2022PRA,RenouBeigi2022PRL} prove rigidity of token-counting distributions on every ring scenario with bipartite sources; their statements concern integer token counts rather than parity, so they do not contain the lemma below, but they are the prior rigidity results on cycles. We give the two proofs because we need the statements for arbitrary latent alphabets and stochastic responses, and because the constant of Lemma~\ref{lem:quantrigidity} enters the distance bounds below.

\begin{lemma}[Quantitative parity rigidity]\label{lem:quantrigidity}
Let $R\in\Ctri$ and $W=\E_R[ABC]$. Then
\[
\max\{\E_RA,\ \E_RB,\ \E_RC\}\ \ge\ -4(1-W).
\]
\end{lemma}

\begin{proof}
Let $a(x,z)$, $b(x,y)$, $c(z,y)$ be the conditional means of $A$, $B$, $C$ given the sources; they take values in $[-1,1]$, and since the three responses use independent private randomness, $W=\E[abc]$ and $\E_RA=\E[a]$, with the expectations over the independent sources $X,Y,Z$. We use two inequalities: for $p,r\in[-1,1]$,
\begin{equation}\label{eq:pr}
|p-r|\le1-pr,
\end{equation}
since for $p\ge r$ the difference $(1-pr)-(p-r)=(1-p)(1+r)$ is nonnegative, and for $u,w\in[-1,1]$,
\begin{equation}\label{eq:uw}
|u+w|\ge2uw,
\end{equation}
since for $uw\le0$ the right side is nonpositive and for $u,w$ of the same sign $|u|+|w|-2|u||w|=|u|(1-|w|)+|w|(1-|u|)\ge0$.

Choose $y_0$ with $\E_{X,Z}[1-a(X,Z)b(X,y_0)c(Z,y_0)]\le1-W$; it exists because the average over $Y$ of the left side is $1-W$. Put $S(X)=b(X,y_0)$ and $T(Z)=c(Z,y_0)$. By \eqref{eq:pr} with $p=a$, $r=ST$,
\[
\E|a-ST|\le\E[1-aST]\le1-W,\qquad\text{so}\qquad |\E[a]-\E[S]\,\E[T]|\le1-W,
\]
using the independence of $X$ and $Z$. Next let $u(y)=\E_X[S(X)b(X,y)]$ and $w(y)=\E_Z[T(Z)c(Z,y)]$, both in $[-1,1]$, and $\upsilon(y)=\sgn(u(y)+w(y))$. By \eqref{eq:uw},
\begin{multline*}
\E_Y[1-\upsilon u]+\E_Y[1-\upsilon w]=\E_Y[2-|u+w|]\le2-2\E_Y[uw]\\
=2-2\E[STbc]\le2-2\bigl(W-\E|a-ST|\bigr)\le4(1-W),
\end{multline*}
where $\E_Y[uw]=\E[S(X)T(Z)b(X,Y)c(Z,Y)]$ by independence of the sources and $\E[STbc]\ge\E[abc]-\E|ST-a|$ because $|bc|\le1$. Hence $\E_Y[1-\upsilon u]\le4(1-W)$ and $\E_Y[1-\upsilon w]\le4(1-W)$. By \eqref{eq:pr} with $p=b(X,Y)$ and $r=S(X)\upsilon(Y)$,
\[
\E|b-S\upsilon|\le\E[1-bS\upsilon]=\E_Y[1-\upsilon u]\le4(1-W),\qquad\text{so}\qquad|\E[b]-\E[S]\,\E[\upsilon]|\le4(1-W),
\]
and likewise $|\E[c]-\E[T]\,\E[\upsilon]|\le4(1-W)$. Write $\sigma=\E[S]$, $\tau=\E[T]$, $\nu=\E[\upsilon]$. If all three means were below $-4(1-W)$, then $\sigma\tau\le\E[a]+(1-W)<0$, $\sigma\nu<0$ and $\tau\nu<0$, whose product $(\sigma\tau\nu)^2$ is nonnegative, a contradiction.
\end{proof}
For $W=1$ the lemma says that some mean is nonnegative, which is Lemma~\ref{lem:parityrigidity} again. A compatible law at total-variation distance $\delta$ from a parity-perfect law has $1-W\le2\delta$, so the lemma also replaces the repair step of \cite{BoreiriEtAl2025}: we construct no repaired model, and the bound holds for arbitrary latent alphabets and stochastic responses.

\begin{lemma}[Distance of the cycle family]\label{lem:cycledistance}
Let $m\ge3$ and let $q>0$ be such that $P_{m,q}$ is a probability law. Then $d_{\TV}(P_{m,q},\Ccyc)\ge q/10$; in particular $P_{m,q}\notin\Ccyc$.
\end{lemma}

\begin{proof}
Partition $\{0,\dots,m-1\}$ into three nonempty contiguous arcs $\Lambda_1,\Lambda_2,\Lambda_3$ and let $\varphi(o)=(V_1,V_2,V_3)$ with $V_r=\prod_{v\in\Lambda_r}o_v$. If $R\in\Ccyc$, then $\varphi_*R\in\Ctri$: the three sources on the arc boundaries are independent, each arc product is a function of the two boundary sources adjacent to its arc together with the sources internal to that arc, and those internal sources can be absorbed as private randomness of the arc. Each $\partial\Lambda_r$ has two elements, and $\partial(\Lambda_1\cup\Lambda_2)=\partial\Lambda_3$, so \eqref{eq:cycletarget} gives
\[
\varphi_*P_{m,q}=\Pi(-q,-q,-q).
\]

Let $R\in\Ccyc$ and $\delta=d_{\TV}(P_{m,q},R)$. Total variation does not increase under $\varphi$, so $Q=\varphi_*R\in\Ctri$ satisfies $d_{\TV}(\Pi(-q,-q,-q),Q)\le\delta$. Since $|\E_Pf-\E_Qf|\le2d_{\TV}(P,Q)$ for $|f|\le1$, the parity $W=\E_Q[V_1V_2V_3]$ satisfies $1-W\le2\delta$ and each mean $\E_QV_r$ is at most $-q+2\delta$. Lemma~\ref{lem:quantrigidity} gives $-q+2\delta\ge-4(1-W)\ge-8\delta$, that is $\delta\ge q/10$.
\end{proof}

\begin{theorem}[Cycles at every order]\label{thm:cycle}
Let $m\ge3$ and $t\ge1$, put
\[
q=\frac1{4m^2t^2},\qquad \eta=\frac q{32m},
\]
and let $\widetilde P_{m,q}=T_\eta^{\otimes m}P_{m,q}$, where $T_\eta$ flips a sign with probability $\eta$. Then $\widetilde P_{m,q}$ is a strictly positive rational law with
\[
\widetilde P_{m,q}\in\Iexp_t(C_m)=\IAI_t(C_m),\qquad
d_{\TV}(\widetilde P_{m,q},\Ccyc)\ \ge\ \frac{11q}{160}=\frac{11}{640\,m^2t^2}>0 .
\]
Consequently no finite order of any of the three hierarchies characterizes $\Ccyc$.
\end{theorem}

\begin{proof}
Apply $T_\eta$ independently to every copied observation of the witness of Lemma~\ref{lem:cyclewitness}. The resulting law is strictly positive, still invariant under the copy-index permutations, and its characters are those of \eqref{eq:masterchar} multiplied by $(1-2\eta)^{|F|}$. The same factor multiplies the target characters, and $|F|$ adds over disjoint blocks, so every injectable marginal and every ancestrally independent product of Lemma~\ref{lem:cyclewitness} persists with $\widetilde P_{m,q}$ in place of $P_{m,q}$, as does the diagonal law. Lemma~\ref{lem:rootsink} again supplies the recursive prescriptions.

For the distance, $d_{\TV}(\widetilde P_{m,q},P_{m,q})\le m\eta=q/32$ by a union bound over the flipped coordinates, and Lemma~\ref{lem:cycledistance} gives $d_{\TV}(\widetilde P_{m,q},\Ccyc)\ge q/10-q/32=11q/160$. Since $t$ was arbitrary and $\widetilde P_{m,q}\notin\Ccyc$, the feasible set at order $t$ is strictly larger than $\Ccyc$ for every $t$.
\end{proof}

For $m=3$ this gives a worst-case distance of order $t^{-2}$ for the triangle, in place of the order $t^{-3}$ that the family of Section~\ref{sec:main} supplies for Question~\ref{q:hausdorff}.

\subsection{A corrected density and a linear lower bound}\label{sec:corrected}

The density \eqref{eq:Hdensity} fixes every auxiliary character, including characters that no prescription refers to. Changing those characters restores positivity at a larger $q$.

\begin{lemma}[Phase telescoping]\label{lem:telescope}
For real $\theta_1,\dots,\theta_m$,
\[
1-\cos\Bigl(\sum_{g=1}^m\theta_g\Bigr)\ \le\ m\sum_{g=1}^m\bigl(1-\cos\theta_g\bigr).
\]
\end{lemma}

\begin{proof}
$1-\cos\phi=\tfrac12|1-e^{i\phi}|^2$, and
$1-\prod_ge^{i\theta_g}=\sum_g\bigl(\prod_{g'<g}e^{i\theta_{g'}}\bigr)(1-e^{i\theta_g})$,
so $|1-\prod_ge^{i\theta_g}|\le\sum_g|1-e^{i\theta_g}|$. Cauchy--Schwarz gives $\bigl(\sum_g a_g\bigr)^2\le m\sum_ga_g^2$.
\end{proof}

\begin{lemma}[Corrected density]\label{lem:corrected}
Fix $m\in\{3,4\}$, $t\ge1$ and $0<q\le1/(16t)$, and set $c=4$ for $m=3$ and $c=5$ for $m=4$. On $m$ families of $t$ signs $s_{g,i}$ put
\begin{equation}\label{eq:Wdensity}
f_g=\prod_{i=1}^t\bigl(1+i\sqrt q\,s_{g,i}\bigr),\qquad R=(1+q)^{t/2},\qquad
W(s)=\Real\prod_{g=1}^mf_g+c\sum_{g=1}^m\bigl(1-\Real f_g\bigr).
\end{equation}
Then $W$ is a real polynomial in $q$ with rational coefficients, $W\ge3/5$ for $m=3$ and $W\ge1/3$ for $m=4$, the law $\rho=2^{-mt}W$ is a probability law, and its characters are
\begin{equation}\label{eq:Wmoments}
\E_\rho\Bigl[\prod_{w\in S}s_w\Bigr]=
\begin{cases}
0,&|S|\text{ odd},\\
(1-c)(-q)^{|S|/2},&S\ne\varnothing\text{ even, inside one family},\\
(-q)^{|S|/2},&|S|\text{ even, meeting at least two families}.
\end{cases}
\end{equation}
\end{lemma}

\begin{proof}
Taking real parts keeps only even powers of $i\sqrt q$, so $W$ is a real polynomial in $q$. Since $|1+i\sqrt q\,s|=(1+q)^{1/2}$ we may write $f_g=Re^{i\theta_g}$, and then
\[
\begin{aligned}
W&=R^m\cos\Bigl(\sum_g\theta_g\Bigr)+c\sum_g\bigl(1-R\cos\theta_g\bigr)\\
&=R^m-cm(R-1)-R^m\Bigl(1-\cos\sum_g\theta_g\Bigr)+cR\sum_g\bigl(1-\cos\theta_g\bigr).
\end{aligned}
\]
Lemma~\ref{lem:telescope} gives
\[
W\ \ge\ R^m-cm(R-1)+\bigl(cR-mR^m\bigr)\sum_g\bigl(1-\cos\theta_g\bigr).
\]
From $tq\le1/16$ we get $R^2=(1+q)^t\le\sum_{j\ge0}(tq)^j\le16/15$ and $R-1=(R^2-1)/(R+1)\le1/30$. For $m=3$, $c=4$: $3R^3\le4R$ because $3R^2\le16/5\le4$, and $W\ge1-12/30=3/5$. For $m=4$, $c=5$: $4R^4\le5R$ because $4R^3\le4(16/15)^{3/2}<5$, and $W\ge1-20/30=1/3$.

For the characters, work under the uniform law on the sign cube, where the families are independent and $\E f_g=1$. For $S$ with parts $S_g$ in the families,
\[
\E\Bigl[\prod_gf_g\prod_{w\in S}s_w\Bigr]=\prod_g\E\Bigl[f_g\prod_{w\in S_g}s_w\Bigr]=(i\sqrt q)^{|S|},
\]
whose real part is $0$ for odd $|S|$ and $(-q)^{|S|/2}$ for even $|S|$. Also $\E[(1-\Real f_g)\prod_{w\in S}s_w]$ vanishes unless $S$ is a nonempty subset of family $g$, in which case it equals $-\Real(i\sqrt q)^{|S|}$. Adding the contributions gives \eqref{eq:Wmoments}; the case $S=\varnothing$ gives $\E_\rho1=1$.
\end{proof}

\begin{theorem}[The square at $q=\Theta(1/t)$]\label{thm:squarelinear}
For every $t\ge1$ and every $0<q\le1/(16t)$, the square parity law $P_q$ of \eqref{eq:squaretarget} lies in $\IAI_t(\square)=\Iexp_t(\square)$.
\end{theorem}

\begin{proof}
Take $m=4$ and $c=5$ in Lemma~\ref{lem:corrected}, name the four families $x,y,z,w$ after the sources $X,Y,Z,W$, and define every copied observation by
\[
A^{il}=x_iw_l,\qquad B^{ij}=x_iy_j,\qquad C^{jk}=y_jz_k,\qquad D^{kl}=z_kw_l,
\]
that is by \eqref{eq:cyclewitness} for $m=4$. The density \eqref{eq:Wdensity} is symmetric within each family, so the law is invariant under independent permutations of the copy indices of each source. By Lemma~\ref{lem:boundary} every boundary occurring in an AI prescription, and every boundary of a subset of one of its blocks, is empty or meets at least two families; by \eqref{eq:Wmoments} such characters take the value $(-q)^{|\partial\cdot|/2}$, unchanged from \eqref{eq:masterchar}. The computation of Lemma~\ref{lem:cyclewitness} therefore applies verbatim: every injectable set has the marginal of $P_{4,q}=P_q$, ancestrally independent injectable blocks are independent, and the diagonal law is $P_q^{\otimes t}$. The moments that Lemma~\ref{lem:corrected} changed are the nonempty even characters inside a single family, and no prescription refers to one. Lemma~\ref{lem:rootsink} gives the recursive prescriptions.
\end{proof}

\begin{theorem}[The triangle at $q=\Theta(1/t)$]\label{thm:trianglelinear}
For every $t\ge1$ and every $0<q\le1/(16t)$, the parity-even law $\Pi(-q,-q,-q)$ of \eqref{eq:paritylaw} lies in $\IAI_t(\triangle)=\Iexp_t(\triangle)$.
\end{theorem}

\begin{proof}
Take $m=3$ and $c=4$ in Lemma~\ref{lem:corrected}, name the three families $x,y,z$ after the sources $X,Y,Z$, and define
\[
A^{ij}=x_iz_j,\qquad B^{ik}=x_iy_k,\qquad C^{jk}=z_jy_k,
\]
which is \eqref{eq:cyclewitness} for $m=3$ in the index convention of Section~\ref{sec:setup}. Positivity and normalization are Lemma~\ref{lem:corrected} with $W\ge3/5$. Symmetry within each family is again immediate. By Lemma~\ref{lem:boundary} every character prescribed by an injectable marginal, by an ancestrally independent product, or by the diagonal law has boundary empty or meeting at least two families, so \eqref{eq:Wmoments} returns $(-q)^{|\partial\cdot|/2}$ and the argument of Lemma~\ref{lem:cyclewitness} applies unchanged. The original copied triangle has $\E A=\E[x_iz_j]=-q$ and likewise $\E B=\E C=-q$, while $ABC=x_iz_jx_iy_kz_jy_k=1$, so its law is $\Pi(-q,-q,-q)=P_{3,q}$. Lemma~\ref{lem:rootsink} gives the recursive prescriptions.
\end{proof}
The exact certificates described in Section~\ref{sec:repro} check the square and triangle witnesses of the two theorems, with and without the flips of Corollaries~\ref{cor:squaredistance} and~\ref{cor:triangledistance}, at orders one and two for the square and one to three for the triangle: positivity of every atom, symmetry, the full diagonal law and every ancestral-independence prescription.


\subsection{Distance and the worst-case profile}\label{sec:worstcase}

For $H\in\{\mathrm{NW},\mathrm{AI},\mathrm{exp}\}$ and a pair-source scenario $G$ set
\begin{equation}\label{eq:Hprofile}
H^H_t(G)=\sup_{P\in\mathcal I^H_t(G)}d_{\TV}(P,\Cscen{G}),
\end{equation}
the worst-case total-variation distance from compatibility of a law the order-$t$ test accepts, with $d_{\TV}$ half the $\ell^1$ distance as in Section~\ref{sec:setup}.

\begin{lemma}[Max-moment inequality for the square]\label{lem:maxmoment}
Every $P\in\Csq$ in sign notation satisfies
\[
\max\{\E A,\ \E B,\ \E[AB]\}+4\,P(ABCD=-1)\ \ge\ 0 .
\]
\end{lemma}

\begin{proof}
Absorb the response seeds into incident sources, so $A=A(X,W)$, $B=B(X,Y)$, $C=C(Y,Z)$, $D=D(Z,W)$ are deterministic signs and $X,Y,Z,W$ are independent. Fix $z$ and write $g_z(Y)=C(Y,z)$, $h_z(W)=D(z,W)$ and
\[
\delta_z=\Pr\bigl(A(X,W)B(X,Y)\ne g_z(Y)h_z(W)\bigr),
\]
so that $\E_Z\delta_Z=P(ABCD=-1)=:\delta$. Put
\[
u_z(x)=\E_W\bigl[A(x,W)h_z(W)\bigr],\qquad v_z(x)=\E_Y\bigl[B(x,Y)g_z(Y)\bigr].
\]
Since $Y$ and $W$ are independent, $\E[ABg_zh_z\mid X=x]=u_z(x)v_z(x)$, and as $ABg_zh_z$ is a sign, $\E_X[1-u_zv_z]=2\delta_z$.

Let $f_z(x)=\sgn(u_z(x)+v_z(x))$, with $\sgn0=1$, and compare $A$ and $B$ with $A'_z=f_z(X)h_z(W)$ and $B'_z=f_z(X)g_z(Y)$. Then
\[
e_z:=\Pr(A\ne A'_z)+\Pr(B\ne B'_z)=\E_X\Bigl[1-\tfrac12|u_z+v_z|\Bigr]\le\E_X\bigl[1-u_zv_z\bigr]=2\delta_z,
\]
using $|u+v|\ge2uv$ on $[-1,1]^2$. Independence of $X,Y,W$ gives
\[
(\E A'_z)(\E B'_z)(\E[A'_zB'_z])=(\E f_z)^2(\E g_z)^2(\E h_z)^2\ \ge\ 0,
\]
so at least one of the three comparison moments is nonnegative. Each differs from the corresponding moment of $P$ by at most $2e_z$, since $|\E U-\E V|\le2\Pr(U\ne V)$ for signs and $\Pr(AB\ne A'_zB'_z)\le\Pr(A\ne A'_z)+\Pr(B\ne B'_z)$. Hence
\[
\max\{\E A,\E B,\E[AB]\}\ \ge\ -2e_z\ \ge\ -4\delta_z .
\]
The left side does not depend on $z$; averaging over $z$ finishes the proof.
\end{proof}

\begin{corollary}[Full-support survivors on the square]\label{cor:squaredistance}
$d_{\TV}(P_q,\Csq)\ge q/6$ for $0<q\le3-2\sqrt2$. With flip probability $\eta=q/64$ applied to each of the four bits and to every copied observation, $\widetilde P_q=T_\eta^{\otimes4}P_q$ is strictly positive, lies in $\IAI_t(\square)=\Iexp_t(\square)$ for $q\le1/(16t)$, and satisfies $d_{\TV}(\widetilde P_q,\Csq)\ge5q/48$. At $q_t=1/(16t)$,
\[
H^{\mathrm{exp}}_t(\square)=H^{\mathrm{AI}}_t(\square)\ \ge\ \frac5{768\,t}.
\]
\end{corollary}

\begin{proof}
If $R\in\Csq$ and $d=d_{\TV}(P_q,R)$, each of $\E_RA$, $\E_RB$, $\E_R[AB]$ is at most $-q+2d$ and the parity error of $R$ is at most $d$, since $P_q$ has none. Lemma~\ref{lem:maxmoment} gives $0\le-q+2d+4d$, so $d\ge q/6$. The noise argument of Theorem~\ref{thm:cycle} applies to the witness of Theorem~\ref{thm:squarelinear}, and $d_{\TV}(\widetilde P_q,P_q)\le4\eta=q/16$, so $d_{\TV}(\widetilde P_q,\Csq)\ge q/6-q/16=5q/48$. Substituting $q_t=1/(16t)$ gives $5/(768t)$, and Lemma~\ref{lem:rootsink} identifies the two suprema.
\end{proof}

\begin{corollary}[Full-support survivors on the triangle]\label{cor:triangledistance}
$d_{\TV}(\Pi(-q,-q,-q),\Ctri)\ge q/10$ for $0<q\le1/3$. With flip probability $\eta=q/192$ on each of the three bits and on every copied observation, the noisy target is strictly positive, lies in $\IAI_t(\triangle)=\Iexp_t(\triangle)$ for $q\le1/(16t)$, and is at distance at least $27q/320$ from $\Ctri$. At $q_t=1/(16t)$,
\[
H^{\mathrm{exp}}_t(\triangle)=H^{\mathrm{AI}}_t(\triangle)\ \ge\ \frac{27}{5120\,t}.
\]
\end{corollary}

\begin{proof}
The distance bound is Lemma~\ref{lem:cycledistance} for $m=3$, whose proof uses only Lemma~\ref{lem:quantrigidity} on the triangle itself. The three flips cost at most $3\eta=q/64$ in total variation, and $q/10-q/64=27q/320$. Feasibility is Theorem~\ref{thm:trianglelinear} with the noise argument of Theorem~\ref{thm:cycle}. Substituting $q_t=1/(16t)$ gives $27/(5120t)$.
\end{proof}

\subsection{A convex-order upper bound}\label{sec:convexorder}

\begin{theorem}[Convex order along the diagonal]\label{thm:convexorder}
Let $G$ be a correlation scenario, that is, finitely many independent latent sources and observed variables with finite alphabets, each a stochastic function of the sources it reads, with $K$ joint outcomes; the order-$n$ inflation and the set $\INW_n(G)$ are defined as in Section~\ref{sec:pairsource} with sources in place of edges. Let $P\in\INW_n(G)$ with witness $\Gamma$. For a deterministic table $\omega$ of copied observations let $q_\omega$ be the observed law obtained by sampling one copy index uniformly and independently for each source and reading $\omega$ at those indices. Then $q_\omega\in\Cscen{G}$ for every $\omega$, and for every convex $\Phi$ on the observed simplex
\begin{equation}\label{eq:convexorder}
\E_\Gamma\,\Phi(q_\omega)\ \le\ \E\,\Phi(\widehat P_n),
\end{equation}
where $\widehat P_n$ is the empirical law of $n$ independent draws from $P$. In particular
\begin{equation}\label{eq:convexl2}
\E_\Gamma\|q_\omega-P\|_2^2\ \le\ \frac{1-\|P\|_2^2}{n},
\qquad
H^{\mathrm{NW}}_n(G)\ \le\ \min\Bigl\{1,\ \frac{\sqrt{K-1}}{2\sqrt n}\Bigr\}.
\end{equation}
\end{theorem}

\begin{proof}
Fix $\omega$. Sampling an index $I_e\in[n]$ uniformly and independently for each source $e$ and letting each observed variable $v$ output $\omega_v\bigl((I_e)_{e\ni v}\bigr)$, where $e\ni v$ ranges over the sources $v$ reads, is a model of $G$ with independent sources $I_e$ and deterministic responses, so $q_\omega\in\Cscen{G}$.

Draw one uniform permutation $\pi_e$ of $[n]$ for each source, independently of each other and of $\omega\sim\Gamma$, and set
\[
T_\ell=\Bigl(\omega_v\bigl((\pi_e(\ell))_{e\ni v}\bigr)\Bigr)_{v},\qquad \ell\in[n],
\qquad \widehat P=\frac1n\sum_{\ell=1}^n\delta_{T_\ell}.
\]
For each fixed choice of permutations, the symmetry of $\Gamma$ carries $(T_1,\dots,T_n)$ to the diagonal rows, so $(T_1,\dots,T_n)\sim P^{\otimes n}$; this remains true after averaging over the permutations, so $\widehat P$ has the law of $\widehat P_n$. Conditionally on $\omega$ and on a fixed $\ell$, the indices $(\pi_e(\ell))_e$ are independent and uniform on $[n]$, so $\E[\delta_{T_\ell}\mid\omega]=q_\omega$ and hence $\E[\widehat P\mid\omega]=q_\omega$. Conditional Jensen gives $\Phi(q_\omega)\le\E[\Phi(\widehat P)\mid\omega]$, and averaging over $\omega$ gives \eqref{eq:convexorder}.

Take $\Phi(\mu)=\|\mu-P\|_2^2$. Then $\E\|\widehat P_n-P\|_2^2=\frac1n\sum_xP(x)(1-P(x))=(1-\|P\|_2^2)/n$, which is the first part of \eqref{eq:convexl2}. Some $\omega$ attains at most the mean, and its $q_\omega$ is compatible, so
\[
d_{\TV}(P,\Cscen{G})\le\tfrac12\|P-q_\omega\|_1\le\frac{\sqrt K}2\|P-q_\omega\|_2\le\frac{\sqrt K}2\sqrt{\frac{1-\|P\|_2^2}n}\le\frac{\sqrt{K-1}}{2\sqrt n},
\]
using $\|P\|_2^2\ge1/K$. Total variation never exceeds $1$.
\end{proof}

The sampling mechanism is that of the finite de Finetti estimate of \cite[Section 4.2 and eq.~(A7)]{NW2020}, which bounds the same quantity by $L(1-\|P\|_2^2)/n$ for $L$ sources through the collision probability of two index samples. Using the full diagonal law through a conditional expectation, rather than its first two degrees, removes the factor $L$. Section~\ref{sec:general} draws the rejecting-order consequence.

\begin{corollary}[Brackets for the square and the triangle]\label{cor:brackets}
For every $t\ge1$,
\[
\frac5{768\,t}\ \le\ H^{\mathrm{exp}}_t(\square)=H^{\mathrm{AI}}_t(\square)\ \le\ H^{\mathrm{NW}}_t(\square)\ \le\ \min\Bigl\{1,\frac{\sqrt{15}}{2\sqrt t}\Bigr\},
\]
\[
\frac{27}{5120\,t}\ \le\ H^{\mathrm{exp}}_t(\triangle)=H^{\mathrm{AI}}_t(\triangle)\ \le\ H^{\mathrm{NW}}_t(\triangle)\ \le\ \frac{\sqrt7}{2\sqrt t}.
\]
\end{corollary}

\begin{proof}
Combine Corollaries~\ref{cor:squaredistance} and~\ref{cor:triangledistance} with Theorem~\ref{thm:convexorder} at $K=16$ and $K=8$, and with Proposition~\ref{prop:gsound}.
\end{proof}

The two exponents do not meet. Whether $H^{\mathrm{NW}}_t(\square)=O(1/t)$ is open, and so is the corresponding question for the triangle; this is the form Question~\ref{q:hausdorff} takes once the lower bounds above replace the family bound.

\subsection{Order conversion on the square}\label{sec:conversion}

\begin{theorem}[NW and AI differ by a bounded change of order]\label{thm:orderconversion}
For every $t\ge2$,
\[
\INW_{\lfloor3t/2\rfloor}(\square)\ \subseteq\ \IAI_t(\square)=\Iexp_t(\square)\ \subseteq\ \INW_t(\square).
\]
\end{theorem}

\begin{proof}
The right inclusion is Proposition~\ref{prop:gsound} and the equality is Lemma~\ref{lem:rootsink}. For the left inclusion let $\Gamma$ be an order-$n$ Navascu\'es--Wolfe witness for $P$ with $n=\lfloor3t/2\rfloor$ and let $\Gamma'$ be its restriction to the copy indices $1,\dots,t$ of each source. Restriction preserves symmetry and the diagonal law of degree $t$, so $\Gamma'\in\INW_t$; we check the AI prescriptions.

Fix an AI set of $\Gamma'$ and decompose it into the connected components of its shared-parent graph, which are injectable and pairwise ancestrally independent. A component's ancestry uses two, three or four source types; the two-type components are single copied observations and correspond to edges of the four-cycle $X-Y-Z-W-X$ of source types. Any two subsets of a four-element set of sizes at least three, or of sizes three and two, intersect, so a component using at least three source types conflicts with every other component and needs a diagonal row of its own. Let $b$ be the number of these.

The remaining components form a multigraph on the source-type four-cycle, with multiplicities $n_A,n_B,n_C,n_D$ on the edges $A,B,C,D$. Opposite edges are disjoint, so the colours $\max(n_A,n_C)$ for the pair $\{A,C\}$ and $\max(n_B,n_D)$ for the pair $\{B,D\}$ suffice, and their sum equals the maximum source-type degree $\Delta$: if $n_A\ge n_C$ and $n_B\ge n_D$ the sum is $n_A+n_B=\deg X$, and $\deg X$ dominates the other three degrees. So $b+\Delta$ rows suffice.

Choose a source type of residual degree $\Delta$, say $X$, together with its two neighbours $Y$ and $W$ on the cycle. A component using at least three source types omits at most one, so it uses at least two of $X,Y,W$; each of the $\Delta$ two-type components incident to $X$ uses exactly two of them. The components have pairwise disjoint copied ancestries, and the three families $X,Y,W$ supply $3t$ copies in $\Gamma'$, so $2b+2\Delta\le3t$ and therefore $b+\Delta\le\lfloor3t/2\rfloor=n$.

Assign each component to the diagonal row of its colour. Components sharing a row use disjoint source types, and distinct components use distinct copies within each family, so independent permutations of the copy indices of the four sources place all components on their assigned rows simultaneously. Under $\Gamma$ the rows are independent with law $P$. Within one row the components have disjoint original source types, so Lemma~\ref{lem:sourcedisjoint}, available because $n\ge2$, makes $P$ factor across them. Hence the joint law of the placed components is the product of their $P$-marginals. Undoing the permutations, the AI set of $\Gamma'$ has exactly the prescribed product.
\end{proof}

The factor $3/2$ is the cost of placing all components on diagonal rows at once; ancestrally disjoint collections of $\lfloor3t/2\rfloor$ injectable components whose source-type conflict graph is complete show that this method cannot do better. Hence $H^{\mathrm{NW}}_{\lfloor3t/2\rfloor}(\square)\le H^{\mathrm{AI}}_t(\square)\le H^{\mathrm{NW}}_t(\square)$, so the two hierarchies have the same polynomial decay exponent if either has one.

\subsection{Low-order thresholds}\label{sec:lowthresholds}

On the parity family we settle the first ten orders exactly. The reduction below turns feasibility for a parity-perfect target into a finite linear program in the numbers of negative auxiliary signs; we state it for the square, where it also covers the Navascu\'es--Wolfe test.

\begin{lemma}[Count-moment reduction]\label{lem:countmoment}
Let $0\le q\le3-2\sqrt2$ and $t\ge1$. Then $P_q\in\IAI_t(\square)$ if and only if there are $w_a\ge0$, $a=(a_X,a_Y,a_Z,a_W)\in\{0,\dots,t\}^4$, with $\sum_aw_a=1$ and
\begin{equation}\label{eq:countmoment}
\sum_aw_a\prod_{g}K^{(t)}_{k_g}(a_g)=(-q)^{S/2},\qquad S=\sum_gk_g,
\end{equation}
for exactly those degree tuples $k\in\{0,\dots,t\}^4$ with $S$ even and $2\max_gk_g\le S$, where
\[
K^{(t)}_k(a)=\binom tk^{-1}\sum_j(-1)^j\binom aj\binom{t-a}{k-j}
\]
is the normalized Krawtchouk polynomial.
\end{lemma}

\begin{proof}[Proof sketch]
Every copied square is injectable, so under any witness it has law $P_q$ and hence parity $+1$ almost surely. Parity forcing across the copy indices then gives sign potentials: from $A^{il}B^{ij}C^{jk}D^{kl}=1$ for all indices, the product $B^{ij}C^{jk}$ depends only on $(i,k)$, so writing $M^{ik}$ for it and $x_i=M^{i1}$, $y_j=C^{j1}$ gives $B^{ij}=x_iy_j$ and, by the same argument on the other pairs, every supported array is of the form $A^{il}=x_iw_l$, $B^{ij}=x_iy_j$, $C^{jk}=y_jz_k$, $D^{kl}=z_kw_l$. The potential is unique up to the simultaneous flip of all four families, so averaging over the two gauges and over independent permutations of the copy indices reduces a witness to a distribution on the four counts $a_g$ of negative signs. Under that symmetrization the average of a character with $k_g$ distinct indices in family $g$ is $\prod_gK^{(t)}_{k_g}(a_g)$, by the definition of the Krawtchouk polynomial as the mean of a $k$-element sign product drawn without replacement.

The degrees that the AI prescriptions constrain are exactly those in the statement. Necessity: by Lemma~\ref{lem:boundary} a nonempty prescribed boundary meets at least two source families, so no single family can carry more than half of the total degree, and the total degree is even. Sufficiency: a balanced degree tuple can be realized by ancestrally disjoint paths in the source four-cycle, matching opposite-type terminals through intermediate copies and then adjacent terminals directly, so the corresponding character is prescribed. The prescribed value is $(-q)^{S/2}$ by \eqref{eq:cycletarget}.
\end{proof}

\begin{lemma}[The Navascu\'es--Wolfe form]\label{lem:countmomentNW}
Under the hypotheses of Lemma~\ref{lem:countmoment}, $P_q\in\INW_t(\square)$ if and only if \eqref{eq:countmoment} holds for exactly the degree tuples in
\[
D_{\mathrm{NW}}(t)=\{b_1+\dots+b_t:\ b_r\in\{0,1\}^4\text{ of even weight}\}\subseteq D_{\mathrm{AI}}(t),
\]
where $D_{\mathrm{AI}}(t)$ is the degree set of Lemma~\ref{lem:countmoment}. The same reductions hold for the triangle family $\Pi(-q,-q,-q)$ with three families of signs and the potentials $A^{ij}=x_iz_j$, $B^{ik}=x_iy_k$, $C^{jk}=z_jy_k$.
\end{lemma}

\begin{proof}[Proof sketch]
Every copied square is carried to the first diagonal row by independent copy-index permutations, so under a symmetric witness with diagonal law $P_q^{\otimes t}$ it has law $P_q$ and parity $+1$; the potential representation and the symmetrization then proceed as in Lemma~\ref{lem:countmoment}. The diagonal law prescribes exactly the characters $\prod_r\prod_{v\in F_r}O_v^{(r)}$ with one vertex set $F_r$ per diagonal row; in potential form row $r$ contributes the indicator of $\partial F_r$ to the degree tuple, and the boundaries of vertex sets of a connected graph are exactly its even-weight edge sets. Conversely a distribution on counts with the prescribed moments yields a symmetric witness with the diagonal law by the construction in the proof of Lemma~\ref{lem:countmoment}. For the triangle, $A^{ij}B^{ik}C^{jk}=1$ at $j=1$ gives $B^{ik}=A^{i1}C^{1k}$, so $x_i=A^{i1}$ and $y_k=C^{1k}$ serve as potentials for $B$, and $z_j=A^{1j}x_1$ completes the representation, unique up to a global flip.
\end{proof}

\begin{lemma}[Monotonicity along the parity direction]\label{lem:paritymonotone}
Let $H\in\{\mathrm{NW},\mathrm{AI},\mathrm{exp}\}$, $G\in\{\triangle,\square\}$, and $0\le q'\le q$ in the range where the parity target is a law. If the parity target with parameter $q$ lies in $I^H_t(G)$, so does the one with parameter $q'$.
\end{lemma}

\begin{proof}
By Lemmas~\ref{lem:countmoment} and~\ref{lem:countmomentNW} we may take a witness to be a law on sign potentials. Multiply every potential sign by an independent $\pm1$ variable of mean $\rho=\sqrt{q'/q}$. A character whose boundary has $S$ elements is multiplied by $\rho^S$, so its value $(-q)^{S/2}$ becomes $(-q')^{S/2}$, while symmetry and the block structure of every prescription are unchanged. Lemma~\ref{lem:rootsink} covers the recursively expressible test.
\end{proof}

By Lemma~\ref{lem:paritymonotone} the set of parameters accepted at a given order is an interval $[0,q^H_t]$, and we bracket $q^H_t$ by a rational primal witness of Lemmas~\ref{lem:countmoment} and~\ref{lem:countmomentNW} at a point below it and rational dual certificates above it.

\begin{proposition}[The parity thresholds at orders up to ten]\label{prop:paritythresholds}
For the square and the triangle, the thresholds $q^{\mathrm{AI}}_t=q^{\mathrm{exp}}_t$ and $q^{\mathrm{NW}}_t$ satisfy the bounds of Table~\ref{tab:thresholds} for $1\le t\le10$, each lower bound being a rational witness and each upper bound a dual certificate valid on the whole interval up to the endpoint of the family. In particular:
\begin{enumerate}[label=(\roman*),nosep]
\item $q^{\mathrm{AI}}_t<q^{\mathrm{NW}}_t$ at every even $t\le10$, so $\IAI_t\subsetneq\INW_t$ there on both scenarios;
\item the two brackets coincide at every odd $t\le9$;
\item $q^{\mathrm{NW}}_{t+1}<q^{\mathrm{AI}}_t$ for every $t\le9$ except $t=3$, where $q^{\mathrm{NW}}_3$ and $q^{\mathrm{NW}}_4$ share the same bracket;
\item on the square, $q^{\mathrm{AI}}_2$ is the zero $r$ of $g(q)=q^3-33q^2+27q-3$ and $P_q\in\INW_2(\square)$ for every $q\le0.17156287$; on the triangle, $q^{\mathrm{AI}}_2=2-\sqrt3$, $q^{\mathrm{NW}}_3=q^{\mathrm{NW}}_4=q^{\mathrm{AI}}_3=1/5$ and $q^{\mathrm{AI}}_5=q^{\mathrm{NW}}_5=1/7$.
\end{enumerate}
Consequently $\INW_3(\square)\subsetneq\IAI_2(\square)=\Iexp_2(\square)\subsetneq\INW_2(\square)$.
\end{proposition}

\begin{table}
\centering
\small
\begin{tabular}{r@{\quad}ll@{\qquad}ll}
\toprule
& \multicolumn{2}{c}{square} & \multicolumn{2}{c}{triangle}\\
$t$ & $q^{\mathrm{AI}}_t$ & $q^{\mathrm{NW}}_t$ & $q^{\mathrm{AI}}_t$ & $q^{\mathrm{NW}}_t$\\
\midrule
1 & $3-2\sqrt2$ & $3-2\sqrt2$ & $1/3$ & $1/3$\\
2 & $0.1324743290$ & $\ge0.1715628752$ & $2-\sqrt3$ & $1/3$\\
3 & $0.0950228950$ & $0.0950228950$ & $1/5$ & $1/5$\\
4 & $0.0835279580$ & $0.0950228950$ & $0.1730047060$ & $1/5$\\
5 & $0.0657380020$ & $0.0657380020$ & $1/7$ & $1/7$\\
6 & $0.0603176600$ & $0.0645356170$ & $0.1283741760$ & $0.1381856900$\\
7 & $0.0502565110$ & $0.0502565110$ & $0.1092774840$ & $0.1092774840$\\
8 & $0.0470271490$ & $0.0494541280$ & $0.1011416960$ & $0.1070388380$\\
9 & $0.0406784540$ & $0.0406784540$ & $0.0888129680$ & $0.0888129680$\\
10 & $0.0384981990$ & $0.0400943600$ & $0.0833179940$ & $0.0870307060$\\
\bottomrule
\end{tabular}
\caption{Certified brackets for the parity thresholds. A decimal entry $x$ means $x\le q^H_t\le x+4\cdot10^{-9}$; the exact entries are the roots of the dual polynomials; the entry $\ge0.1715628752$ is a lower bound only, the family ending at $3-2\sqrt2\approx0.1715728753$.}\label{tab:thresholds}
\end{table}

\begin{proof}
All statements are exact certificates, described in Section~\ref{sec:repro}: for each order and hierarchy a nonnegative rational count witness at the lower endpoint, checked against every prescribed moment of Lemmas~\ref{lem:countmoment} and~\ref{lem:countmomentNW}, and a chain of rational dual functionals, each nonpositive on all count configurations, whose prescribed expectations are polynomials in $q$ certified positive by Sturm sequences on intervals covering everything above the upper endpoint. Lemma~\ref{lem:paritymonotone} turns the point witnesses into the intervals. The exact values arise where the dual polynomial factors: at square order two the dual is the functional $D$ with $\E D=(1+q)g(q)$, verified coefficient by coefficient; on the triangle the duals are $-(q+1)(q^2-4q+1)$ at order two and $(5q-1)(q+1)^2$, $(5q-1)(q+1)^3$ at orders three and four, and the order-five dual vanishes at $1/7$. The degree sets $D_{\mathrm{AI}}(t)$ and $D_{\mathrm{NW}}(t)$ are enumerated by brute force over all ancestral-independence families for $t\le3$ and, for every $t\le10$, every tuple of $D_{\mathrm{AI}}(t)$ is realized by an explicit family that the verifier checks against the definitions. The chain follows from the order-two and order-three brackets on the square, with the equality from Lemma~\ref{lem:rootsink}.
\end{proof}

The table shows two regularities that the theorems do not explain. The products $tq^H_t$ increase towards constants near $0.4$ on the square and $0.9$ on the triangle, an order of magnitude above the $1/16$ of Theorems~\ref{thm:squarelinear} and~\ref{thm:trianglelinear}; and the two hierarchies coincide along this direction at every odd order up to nine while separating at every even order, with a ratio $q^{\mathrm{NW}}_t/q^{\mathrm{AI}}_t-1$ decreasing like $0.4/t$. Neither pattern is proved beyond the certified orders; uncertified numerics at triangle order fifteen show a relative gap of $10^{-4}$ between the two hierarchies, which may be the end of the coincidence or conditioning error.

\providecommand{\Cpath}{\mathcal C_{P_5}}
\providecommand{\Real}{\operatorname{Re}}

\subsection{The five-observer path}\label{sec:fivepath}

A path is a tree, and trees carry no parity obstruction of the kind used in Section~\ref{sec:cycles}. The scenario, in the language of Section~\ref{sec:pairsource}, is
\[
A(X),\qquad B(X,L),\qquad C(L,R),\qquad D(R,Z),\qquad E(Z),
\]
with $X,L,R,Z$ mutually independent (Figure~\ref{fig:path}). We write $\Cpath$ for its compatible set, record the endpoint observations $A,E$ as bits $x,z$ and the middle observations $B,C,D$ as signs. The order-$t$ inflation of Definition~\ref{def:ginflation} has copies $X_a,L_i,R_j,Z_e$, $a,i,j,e\in[t]$, and the $3t^2+2t$ copied observations
\[
A^a,\qquad B^{ai},\qquad C^{ij},\qquad D^{je},\qquad E^e .
\]

\begin{figure}[htbp]
\centering
\begin{tikzpicture}[x=1cm,y=1cm]
\foreach \v/\x in {A/0,B/2.4,C/4.8,D/7.2,E/9.6}{\node[observer] (\v) at (\x,0) {$\v$};}
\foreach \s/\x/\a/\b in {X/1.2/A/B,L/3.6/B/C,R/6/C/D,Z/8.4/D/E}{
\node[source] (\s) at (\x,1.2) {$\s$};
\draw[causal] (\s)--(\a);\draw[causal] (\s)--(\b);
}
\node[font=\small,below=3mm of A] {$x\in\{0,1\}$};
\node[font=\small,below=3mm of E] {$z\in\{0,1\}$};
\node[font=\small,align=center] at (4.8,-1.25) {Condition on the endpoint outcomes $x,z$:\\$f_{xz}=\mathbb E_{L,R}[b_x(L)c(L,R)d_z(R)]$};
\draw[inkblue,line width=1pt] (2.1,-.6)--(7.5,-.6);
\node[font=\small,align=center] at (4.8,-2.25) {Compatibility: $\sqrt{|I|}+\sqrt{|J|}\le1$\qquad Target: $I=J=(1+h)/4$};
\end{tikzpicture}
\caption{The five-observer path. Conditioning on $A=x$ and $E=z$ changes only the outer sources $X,Z$; the inner sources $L,R$ remain independent. This gives the bilocal inequality. The target violates it by $\sqrt{1+h}-1$ while retaining full support.}\label{fig:path}
\end{figure}


\subsubsection{The target}\label{sec:pathtarget}

\begin{definition}[Path target]\label{def:pathtarget}
For $0<h<1$ define the law $P_h$ by
\begin{equation}\label{eq:pathtarget}
P_h(x,b,c,d,z)=\frac1{32}\Bigl[1+bcd\,\frac{1+h}4\bigl(1+(-1)^{x+z}\bigr)\Bigr],
\qquad x,z\in\{0,1\},\ b,c,d\in\{-1,1\}.
\end{equation}
\end{definition}
\FloatBarrier

\begin{lemma}[Properties of $P_h$]\label{lem:pathtarget}
$P_h$ is a probability law, every atom is at least $(1-h)/64$, the pair $(A,E)$ consists of independent fair bits, every proper nonempty moment of $(B,C,D)$ conditional on $(A,E)$ vanishes, and
\begin{equation}\label{eq:fxz}
f_{xz}:=\E\bigl[BCD\mid A=x,E=z\bigr]=\frac{1+h}4\bigl(1+(-1)^{x+z}\bigr)
=\begin{cases}(1+h)/2,&x=z,\\0,&x\ne z.\end{cases}
\end{equation}
\end{lemma}

\begin{proof}
The bracket in \eqref{eq:pathtarget} lies between $1-(1+h)/2=(1-h)/2$ and $1+(1+h)/2$, so every atom lies in $[(1-h)/64,\,(3+h)/64]$ and is positive. Summing $bcd$ over $\{-1,1\}^3$ gives zero, so the atoms sum to $32\cdot\frac1{32}=1$ and the $(A,E)$ marginal is uniform on four cells. Given $(x,z)$, the conditional density of $(b,c,d)$ is $\frac18\bigl[1+bcd\,f_{xz}\bigr]$, whose only nonvanishing nontrivial character is $bcd$, with value $f_{xz}$; that is \eqref{eq:fxz} and the vanishing of the proper moments.
\end{proof}

\subsubsection{A bilocal obstruction}\label{sec:bilocal}

\begin{lemma}[Bilocal inequality on the path]\label{lem:bilocal}
Let $R\in\Cpath$ with $R(A=x,E=z)>0$ for all four $(x,z)$, put $f_{xz}=\E_R[BCD\mid A=x,E=z]$ and
\[
I=\frac14\sum_{x,z}f_{xz},\qquad J=\frac14\sum_{x,z}(-1)^{x+z}f_{xz}.
\]
Then $\sqrt{|I|}+\sqrt{|J|}\le1$.
\end{lemma}

\begin{proof}
Absorb the response seeds into incident sources. Conditioning on $\{A=x\}$ reweights the law of $X$ alone, and conditioning on $\{E=z\}$ reweights the law of $Z$ alone, so under the conditional law $X,L,R,Z$ are still mutually independent, and $L,R$ keep their original laws. Writing
\[
b_x(\ell)=\E\bigl[B\mid L=\ell,\ A=x\bigr],\qquad
d_z(\rho)=\E\bigl[D\mid R=\rho,\ E=z\bigr],
\]
\[
c(\ell,\rho)=\E\bigl[C\mid L=\ell,\ R=\rho\bigr],
\]
all with values in $[-1,1]$, conditional independence of the three responses given the sources gives
\[
f_{xz}=\E_{L,R}\bigl[b_x(L)\,c(L,R)\,d_z(R)\bigr].
\]
Put $\beta_\pm=\tfrac12(b_0\pm b_1)$ and $\delta_\pm=\tfrac12(d_0\pm d_1)$. Then
\[
I=\E_{L,R}\bigl[\beta_+(L)\,c(L,R)\,\delta_+(R)\bigr],\qquad
J=\E_{L,R}\bigl[\beta_-(L)\,c(L,R)\,\delta_-(R)\bigr],
\]
so, using $|c|\le1$ and the independence of $L$ and $R$,
\[
|I|\le p_+r_+,\qquad |J|\le p_-r_-,\qquad
p_\pm=\E_L|\beta_\pm|,\quad r_\pm=\E_R|\delta_\pm| .
\]
For real $u,v$ with $|u|,|v|\le1$,
$\bigl|\tfrac{u+v}2\bigr|+\bigl|\tfrac{u-v}2\bigr|=\max(|u|,|v|)\le1$,
so $p_++p_-\le1$ and $r_++r_-\le1$. Cauchy--Schwarz gives
\[
\sqrt{|I|}+\sqrt{|J|}\le\sqrt{p_+r_+}+\sqrt{p_-r_-}\le\sqrt{p_++p_-}\,\sqrt{r_++r_-}\le1 . \qedhere
\]
\end{proof}

Inequalities of this shape go back to Branciard, Rosset, Gisin and Pironio \cite{BranciardRossetGisinPironio2012} for the bilocal scenario with measurement inputs; Lemma~\ref{lem:bilocal} is the specialization to a single fixed setting, one binary central output and two binary endpoint observations, proved here for arbitrary latent alphabets and stochastic responses.

\begin{corollary}[Quantitative incompatibility]\label{cor:pathdistance}
For every $0<h<1$,
\[
P_h\notin\Cpath\qquad\text{and}\qquad d_{\TV}(P_h,\Cpath)\ \ge\ \frac h{96}.
\]
\end{corollary}

\begin{proof}
By \eqref{eq:fxz}, $I=J=(1+h)/4$ for $P_h$, so $\sqrt I+\sqrt J=\sqrt{1+h}>1$ and Lemma~\ref{lem:bilocal} rules out compatibility.

Let $R\in\Cpath$ and $d=d_{\TV}(P_h,R)$. If $d\ge1/8$ then $d>1/96>h/96$, so assume $d<1/8$. Each endpoint cell satisfies $|R(x,z)-\tfrac14|\le d$, hence $R(x,z)\ge\tfrac14-d>\tfrac18>0$ and Lemma~\ref{lem:bilocal} applies to $R$. Write $N_R(x,z)=\E_R\bigl[BCD\,\one\{A=x,E=z\}\bigr]$ and $N_P$ for the same quantity under $P_h$, so $f^R_{xz}=N_R(x,z)/R(x,z)$ and likewise for $P_h$. The eight atoms of one endpoint cell contribute at most $2d$ to $\sum|R-P_h|$, so $|N_R-N_P|\le2d$, while $|R(x,z)-P_h(x,z)|\le d$. Since $|N_P|\le P_h(x,z)=\tfrac14$,
\[
\bigl|f^R_{xz}-f^P_{xz}\bigr|
\le\frac{|N_R-N_P|}{R(x,z)}+\frac{|N_P|\,\bigl|P_h(x,z)-R(x,z)\bigr|}{R(x,z)P_h(x,z)}
\le\frac{2d+d}{R(x,z)}\le24d .
\]
Averaging, $|I_R-I_P|\le24d$ and $|J_R-J_P|\le24d$. If $24d<h/4$ then
$I_R,J_R>\tfrac{1+h}4-\tfrac h4=\tfrac14$, so $\sqrt{|I_R|}+\sqrt{|J_R|}>1$, contradicting Lemma~\ref{lem:bilocal}. Hence $24d\ge h/4$.
\end{proof}

\subsubsection{The witness at every order}\label{sec:pathwitness}

\begin{theorem}[The path at every order]\label{thm:fivepath}
For every $t\ge1$ put $\gamma=1/(4t)$ and $h_t=\gamma^2=1/(16t^2)$. Then
\[
P_{h_t}\in\Iexp_t(P_5)=\IAI_t(P_5),\qquad P_{h_t}\notin\Cpath ,
\]
with a strictly positive rational target and a rational witness. Consequently $\tmin^H(P_{h_t})>t$ for $H\in\{\mathrm{NW},\mathrm{AI},\mathrm{exp}\}$, no finite order of any of the three hierarchies characterizes $\Cpath$, and
\[
d_{\TV}(P_{h_t},\Cpath)\ \ge\ \frac1{1536\,t^2}.
\]
\end{theorem}

\begin{proof}
\emph{The auxiliary law.} On signs $u=(u_1,\dots,u_t)$ and $v=(v_1,\dots,v_t)$ set
\begin{equation}\label{eq:Ht}
H_t(u,v)=2^{-2t}\,\Real\Bigl[\prod_{i=1}^t(1+i\gamma u_i)\prod_{j=1}^t(1-i\gamma v_j)\Bigr].
\end{equation}
Writing $\sigma=(u_1,\dots,u_t,-v_1,\dots,-v_t)\in\{-1,1\}^{2t}$, the bracket is $\prod_{k=1}^{2t}(1+i\gamma\sigma_k)$, whose real part keeps the even characters with coefficient $(i\gamma)^{|S|}=(-h_t)^{|S|/2}$, so
\[
H_t=2^{-2t}\sum_{\substack{S\subseteq[2t]\\|S|\text{ even}}}(-h_t)^{|S|/2}\prod_{k\in S}\sigma_k .
\]
This is $H_{2t,h_t}$ of Lemma~\ref{lem:parityaux} read in the coordinates $\sigma$, and $(2t)^2h_t=1/4$, so $H_t$ is a probability law with
\begin{equation}\label{eq:Htpos}
H_t(u,v)\ \ge\ \tfrac23\,2^{-2t}\ >\ 0 .
\end{equation}
Its values are rational.

\emph{The witness.} Draw $(u,v)\sim H_t$; independently draw fair bits $X_a$ and $Z_e$, fair signs $r_i$ and $s_j$, and fair signs $\nu_{ij}$, all mutually independent, $a,i,j,e\in[t]$. With $u^0=1$ and $u^1=u$, define every copied observation by
\begin{equation}\label{eq:pathwitness}
A^a=X_a,\qquad
B^{ai}=r_i\,u_i^{X_a},\qquad
C^{ij}=\begin{cases}r_is_j,&u_i=v_j,\\ r_is_j\nu_{ij},&u_i\ne v_j,\end{cases}\qquad
D^{je}=s_j\,v_j^{Z_e},\qquad
E^e=Z_e .
\end{equation}
Each copied observation is a function of the data attached to its copied parents together with its own private sign: $r_i$ and $u_i$ belong to $L_i$, and $s_j$ and $v_j$ to $R_j$. Call the resulting law $\Gamma_t$. It is a nonnegative rational probability law. The auxiliary density is strictly positive, but the observed table can have zero entries. As in Section~\ref{sec:construction}, $\Gamma_t$ is a law on the copied observations and does not come from an inflated model, since under $H_t$ the families $u$ and $v$ are correlated.

\emph{Formal weights.} Give each $u_i$ the complex weight $(1+i\gamma u_i)/2$ and each $v_j$ the weight $(1-i\gamma v_j)/2$, and keep the actual probabilities of $X,Z,r,s,\nu$. Write $\E_\ast$ for the resulting formal expectation, a product functional on the $6t+t^2$ auxiliary coordinates, with $\E_\ast u_i=i\gamma$ and $\E_\ast v_j=-i\gamma$. Since the weights of the auxiliary coordinates multiply to a complex number whose real part is $H_t$, and the other coordinates carry real weights,
\begin{equation}\label{eq:realpart}
\E_{\Gamma_t}[F]=\Real\,\E_\ast[F]\qquad\text{for every real-valued }F .
\end{equation}

\emph{Every copied path has law $P_h$.} Fix $(a,i,j,e)$ and condition on $X_a=x$, $Z_e=z$, which are fair and independent and reproduce the $(A,E)$ marginal of $P_h$. The masks $r_i$ and $s_j$ are fair signs independent of everything else and each occurs in exactly one of $B^{ai}$, $D^{je}$ and in $C^{ij}$, so every product of a proper nonempty subset of $\{B^{ai},C^{ij},D^{je}\}$ carries $r_i$ or $s_j$ to an odd power and has formal expectation zero. In the triple product the masks cancel:
\[
B^{ai}C^{ij}D^{je}=u_i^{\,x}\,v_j^{\,z}\cdot\bigl(\one\{u_i=v_j\}+\nu_{ij}\one\{u_i\ne v_j\}\bigr),
\]
and $\E_\ast\nu_{ij}=0$, so $\E_\ast\bigl[B^{ai}C^{ij}D^{je}\bigr]=\E_\ast\bigl[u_i^{\,x}v_j^{\,z}\one\{u_i=v_j\}\bigr]$. The formal weight of $(u_i,v_j)=(+1,+1)$ is $\frac{1+i\gamma}2\cdot\frac{1-i\gamma}2=\frac{1+h}4$, and the weight of $(-1,-1)$ is the same real number, so
\[
\E_\ast\bigl[u_i^{\,x}v_j^{\,z}\one\{u_i=v_j\}\bigr]=\frac{1+h}4\bigl(1+(-1)^{x+z}\bigr)=f_{xz},
\]
which is real. Every character of the copied path is thus one of: a character of the fair bits $X_a,Z_e$ alone; a character carrying a proper nonempty subset of $\{B^{ai},C^{ij},D^{je}\}$, with formal expectation zero; or a character carrying all three, with formal expectation $\frac14\sum_{x,z}\chi(x,z)f_{xz}$. These are real and agree with the corresponding characters of $P_h$ in Lemma~\ref{lem:pathtarget}, so by \eqref{eq:realpart}
\[
\Law_{\Gamma_t}\bigl(A^a,B^{ai},C^{ij},D^{je},E^e\bigr)=P_h,
\]
and $\E_\ast[G]$ is real for every function $G$ of a copied path.

\emph{Injectable marginals.} An injectable set of the path contains at most one copied observation over each of $A,B,C,D,E$ and its copy indices agree on shared sources, so it lies in a copied path $(a,i,j,e)$ with consistent indices, after completing unused indices arbitrarily. Its marginal is therefore the corresponding marginal of $P_h$, and $\E_\ast[G]$ is real for every function $G$ of it.

\emph{Ancestral products and the diagonal law.} Let $F_1,\dots,F_r$ be pairwise ancestrally independent injectable sets. Distinct copied sources carry distinct auxiliary coordinates and distinct masks, and distinct copied observations carry distinct private signs, so the $F_k$ depend on disjoint coordinate families and $\E_\ast\bigl[\prod_k G_k\bigr]=\prod_k\E_\ast[G_k]$ for functions $G_k$ of $F_k$. Each factor is real by the previous paragraph, hence so is the product, and \eqref{eq:realpart} gives
\[
\E_{\Gamma_t}\Bigl[\prod_kG_k\Bigr]=\prod_k\E_\ast[G_k]=\prod_k\E_{\Gamma_t}[G_k].
\]
Taking the $G_k$ to be indicators of atoms shows the blocks are independent with the prescribed marginals. The diagonal rows $(A^\ell,B^{\ell\ell},C^{\ell\ell},D^{\ell\ell},E^\ell)$, $\ell\in[t]$, are pairwise ancestrally independent injectable copied paths, so the diagonal law is $P_h^{\otimes t}$.

\emph{Symmetry.} The density \eqref{eq:Ht} is symmetric in $u_1,\dots,u_t$ and in $v_1,\dots,v_t$ separately, the families $(X_a)$, $(Z_e)$, $(r_i)$, $(s_j)$ and $(\nu_{ij})$ are i.i.d., and \eqref{eq:pathwitness} is equivariant, so $\Gamma_t$ is invariant under independent permutations of the four copy-index families. Together with the previous two paragraphs this places $P_{h_t}$ in $\IAI_t(P_5)$, and Lemma~\ref{lem:rootsink} gives $\Iexp_t(P_5)=\IAI_t(P_5)$.

\emph{Conclusion.} Corollary~\ref{cor:pathdistance} gives $P_{h_t}\notin\Cpath$ and
$d_{\TV}(P_{h_t},\Cpath)\ge h_t/96=1/(1536t^2)$. Membership at order $t$ forces $\tmin^H(P_{h_t})>t$ for each of the three hierarchies, by Proposition~\ref{prop:gsound}.
\end{proof}

Along this family the first rejecting order grows without bound while the target stays uniformly inside the simplex: since $h_t\le1/16$, Lemma~\ref{lem:pathtarget} gives
\[
P_{h_t}(x,b,c,d,z)\ \ge\ \frac{1-h_t}{64}\ \ge\ \frac{15}{1024}
\]
for every $t$ and every atom. Vanishing probabilities play no role in the obstruction. The theorem gives a distance lower bound of order $t^{-2}$ and leaves the distance asymptotics of this family undetermined. The general upper bound is $O(t^{-1/2})$. It remains open whether the path admits accepted laws at distance $\Omega(t^{-1})$, as the triangle and square do.

